\subsubsection{Resumen ejecutivo}
LogiMarket 360 controla la venta y entrega de accesorios informáticos mediante un catálogo con kits anidados, pedidos con instantáneas, cobro simulado, existencias y devoluciones. Persiste en SQLite los efectos comerciales y los eventos que deben procesarse después de confirmar. La interfaz utiliza Windows Forms y compara los permisos visibles con la autorización efectiva de Application.

El informe documenta la revisión \ruta{12458bc1cf073e2575e1d85bda0f9f0517a1fa19}. La ejecución automatizada archivada acredita 119 pruebas aprobadas. La matriz de auditoría registra los 38 casos de aceptación y los 23 patrones. El resultado global de entrega sigue condicionado por la verificación de presentación al 125\,\% y al 150\,\%; este requisito es transversal y se detalla en la sección~\ref{sec:dpi}.

\subsubsection{Problema, actores y requisitos}
La empresa necesita controlar existencias físicas al vender tanto hojas como kits, conservar el valor realmente acordado y evitar cobros o reembolsos repetidos. El diseño separa la operación comercial, sus efectos confirmados y la notificación recuperable. Los comprobantes, el proveedor de pago, los avisos y la guía de envío son simulados e internos. La fuente normativa es la guía de práctica \cite{ape04}.

\diagrama{contexto}{contexto}{Contexto operativo en un equipo Windows y sus límites.}
\diagrama{casos}{casos}{Casos de uso y responsabilidades principales por rol.}

\begin{longtable}{@{}L{.19\linewidth}L{.76\linewidth}@{}}
\toprule Rol & Operaciones y alcance \\ \midrule\endhead
ADMIN & Todas las operaciones, incluyendo mantenimiento, demostración y auditoría. \\
VENTAS & Gestión de clientes y pedidos, validación, confirmación y cobro. También consulta el catálogo operativo. \\
BODEGA & Consultas de preparación e inventario, despacho y entrega. El Proxy deniega un intento de cobro antes de las trazas y del proveedor. \\
CONSULTA & Operaciones de lectura, incluidos los reportes permitidos. Sin mutaciones comerciales ni ajustes. \\
\bottomrule
\end{longtable}
\noindent Los permisos se verifican con \ruta{AutorizacionOperaciones.ExigirAsync}. Las pruebas con sesiones por rol contrastan las operaciones de escritura y de consulta. La Figura~\ref{diag:casos} representa el alcance principal, que se completa con las consultas autorizadas en Application.

\paragraph{Trazabilidad funcional y de calidad.}
La siguiente matriz vincula cada requisito con escenarios y evidencia. COMPLETO significa que el alcance auditado cuenta con evidencia. PARCIAL indica una condición pendiente. NO VERIFICADO identifica una ejecución no corroborada. INCUMPLIDO se reserva para una condición que la ejecución contradice.
\small
\begin{longtable}{@{}L{.09\linewidth}L{.47\linewidth}L{.35\linewidth}@{}}
\caption{Trazabilidad de requisitos funcionales y no funcionales.}\\
\toprule ID & Condición de la guía & Evidencia y estado \\ \midrule\endfirsthead
\toprule ID & Condición de la guía & Evidencia y estado \\ \midrule\endhead
RF01 & Identificarse y cerrar sesión localmente con usuarios iniciales, hash y cuatro roles. & P01, P22; login correcto/incorrecto y sesión cerrada. \textbf{COMPLETO} \\[2pt]
RF02 & Consultar y administrar productos con CRUD, búsqueda e invariantes del catálogo. & P01, P21; duplicado SKU y cantidades inválidas. \textbf{COMPLETO} \\[2pt]
RF03 & Registrar y editar kits, impedir ciclos y expandir SKU finales. & P05, P18, P27, P34. \textbf{COMPLETO} \\[2pt]
RF04 & Buscar, registrar y editar clientes con identificación, categoría, contacto y estado activo. & P23; ejercicio CRUD y cliente inactivo. \textbf{COMPLETO} \\[2pt]
RF05 & Crear/editar borrador con cliente, líneas, cantidades, entrega e instantáneas; rechazar vacío. & P08, P13; modificar catálogo posterior y verificar instantánea. \textbf{COMPLETO} \\[2pt]
RF06 & Clonar carrito con copia profunda y deshacer/rehacer hasta cinco estados completos, solo en borrador. & P14, P29. \textbf{COMPLETO} \\[2pt]
RF07 & Calcular subtotal, promoción única, flete, plazo y total con cifras verificables. & P02–P04, P12, P32; umbral 9/10 y redondeo. \textbf{COMPLETO} \\[2pt]
RF08 & Validar por cadena, cancelar únicamente en estados permitidos y consultar historial. & P08, P23, P24, P31. \textbf{COMPLETO} \\[2pt]
RF09 & Confirmar con cobro simulado sustituible, registrar trazas y evitar pago doble/sobreventa. & P02–P07, P16, P19–P20, P22, P36. \textbf{COMPLETO} \\[2pt]
RF10 & Preparar, despachar y entregar en orden, generar guía simulada y consultar seguimiento. & P09, P12, P31; trazabilidad de guía/entrega. \textbf{COMPLETO} \\[2pt]
RF11 & Devolver totalmente desde ENTREGADO conforme a motivo, plazo y condición; reintegro y reembolso únicos. & P33, P35; límites inclusivos y producto usado. \textbf{COMPLETO} \\[2pt]
RF12 & Consultar/registrar inventario y movimientos con motivo y referencia, incluido ajuste auditado. & P02–P07, P35–P36; ajuste positivo/negativo. \textbf{COMPLETO} \\[2pt]
RF13 & Obtener reporte por fechas y filtros por estado, total y cliente con AND/OR y paréntesis. & P25, P26, P34; filtro por rango y resumen. \textbf{COMPLETO} \\[2pt]
RF14 & Emitir comprobantes internos de venta y devolución a vista previa o TXT. & P17; cuatro combinaciones. \textbf{COMPLETO} \\[2pt]
RF15 & Generar avisos simulados EMAIL/SMS, auditoría y métricas por pedido, despacho y devolución. & P11, P15, P19, P28, P30, P37. \textbf{COMPLETO} \\[2pt]
RF16 & Mostrar panel con pedidos, ventas, alertas, stock y métricas; refrescar catálogo, carrito y pedidos mediante mediador. & P15, P21, P28, P30. \textbf{COMPLETO} \\[2pt]
RF17 & Demostrar cambio de proveedor, fallo de cobro y fallo/reintento Outbox en entorno controlado. & P06, P16, P37. \textbf{COMPLETO} \\[2pt]
RF18 & Persistir pedido, pago, stock, despacho, devolución y auditoría al reiniciar. & P01, P10, P35, P37; consulta postreinicio. \textbf{COMPLETO} \\[2pt]
RNF01 & Integridad. Cero efectos de negocio parciales en rechazo. Un pago por pedido, un movimiento por SKU en confirmación, una devolución y stock $\ge$ 0. & P05–P08, P20, P35–P36, índices, estado de DB. \textbf{COMPLETO} \\[2pt]
RNF02 & Mantenibilidad. Cinco proyectos y cero dependencias prohibidas. Patrones, SOLID, Repository, CQRS básico, transacción real y DI por constructor. & P38, \_build\_ limpio, .csproj, diagramas, código e invocaciones reales. \textbf{COMPLETO} \\[2pt]
RNF03 & Usabilidad. Doce áreas accesibles desde la ventana principal; TAB, foco, etiquetas, validación próxima al campo, escala, mensajes, confirmación, selección persistente y dos decimales. & Cotejo de las 12 vistas y alcance en la sección~\ref{sec:dpi}. \textbf{PARCIAL} \\[2pt]
RNF04 & Rendimiento. P21 ejecutado sobre al menos 1.000 filas, con identidad compartida, instancias y memoria frente a escenario equivalente. La guía no fija latencia ni porcentaje de ahorro. & Medición reproducible, condiciones y resultado bruto obtenido. \textbf{COMPLETO} \\[2pt]
RNF05 & Seguridad local. Cuatro roles, autorización de operaciones y cero contraseñas en claro; denegación de pago antes de trazas y proveedor. & P01, P22, inspección de BD, revisión del ZIP. \textbf{COMPLETO} \\[2pt]
RNF06 & Trazabilidad. Registrar cambios de estado, movimientos, correlación, comandos, eventos y evidencias para todos los requisitos/patrones y casos. & Matrices, logs, captura, referencia a archivo/método/commit. \textbf{COMPLETO} \\[2pt]
RNF07 & Testabilidad. Ejecutar P01–P38 y aprobar al menos 25 pruebas automáticas en dominio, aplicación e integración. & P38, informe TRX, datos repetibles y resultados obtenidos. \textbf{COMPLETO} \\[2pt]
RNF08 & Recuperabilidad. Base vacía recuperable vía migraciones/seed, persistencia al reiniciar y evento reintentable sin duplicación. & P01, P10, P37, instrucciones reproducibles. \textbf{COMPLETO} \\[2pt]
RNF09 & Compatibilidad y alcance. Windows, Visual Studio 2022 compatible, C\#/.NET 8, WinForms, EF Core 8, SQLite y \ruta{Microsoft.Extensions.DependencyInjection}. Colaboradores simulados y comprobantes internos. & Compilación, UI, esquema relacional real y demostración sin servicios externos. \textbf{COMPLETO} \\[2pt]
RNF10 & Calidad y autoría. Registrar al menos diez riesgos reales con refactorizaciones, Git y pruebas antes/después. Aplicar cinco principios SOLID y Clean Code/KISS/DRY/YAGNI. & Registro de riesgos e intervenciones en la sección~\ref{sec:riesgos}. \textbf{PARCIAL} \\[2pt]
\bottomrule\end{longtable}\normalsize


